% Felipe Muñoz Casasola
% This is free and unencumbered software released into the public domain.

% Anyone is free to copy, modify, publish, use, compile, sell, or
% distribute this software, either in source code form or as a compiled
% binary, for any purpose, commercial or non-commercial, and by any
% means.

% In jurisdictions that recognize copyright laws, the author or authors
% of this software dedicate any and all copyright interest in the
% software to the public domain. We make this dedication for the benefit
% of the public at large and to the detriment of our heirs and
% successors. We intend this dedication to be an overt act of
% relinquishment in perpetuity of all present and future rights to this
% software under copyright law.

% THE SOFTWARE IS PROVIDED "AS IS", WITHOUT WARRANTY OF ANY KIND,
% EXPRESS OR IMPLIED, INCLUDING BUT NOT LIMITED TO THE WARRANTIES OF
% MERCHANTABILITY, FITNESS FOR A PARTICULAR PURPOSE AND NONINFRINGEMENT.
% IN NO EVENT SHALL THE AUTHORS BE LIABLE FOR ANY CLAIM, DAMAGES OR
% OTHER LIABILITY, WHETHER IN AN ACTION OF CONTRACT, TORT OR OTHERWISE,
% ARISING FROM, OUT OF OR IN CONNECTION WITH THE SOFTWARE OR THE USE OR
% OTHER DEALINGS IN THE SOFTWARE.

% For more information, please refer to <https://unlicense.org/>

\documentclass{article}

\usepackage[spanish]{babel}
\usepackage[margin=3cm]{geometry}
\usepackage{microtype}
\usepackage[pdfusetitle]{hyperref}
\usepackage{minted}

\title{Instalación de Forgejo}
\author{Felipe Muñoz Casasola}
\date{11 de marzo de 2026}

\begin{document}
    \maketitle
    \clearpage
    \tableofcontents
    \clearpage

    \section{Muy importante}
    Todo repositorio en mi instancia de Forgejo se deberá configurar el empuje
    automático de repositorios a la instancia de Forgejo que deberé tener instalada
    en el móvil. Por tanto, \textbf{solamente podré iniciar la instancia cuando
    esté conectado al WiFi de móvil por USB}. El detalle del USB es para todavía
    mayor seguridad. De esta forma, la instancia siempre estará totalmente asilada
    de otros dispositivos que no sean ni mi ordenador ni mi móvil. A los que no
    les debe importar lo que yo haga.

    Es posible que tenga que usar \texttt{nmcli} para ver las conexiones disponibles,
    la del móvil comenzará por \texttt{enx}, y entonces me desconecto con \texttt{nmcli device
    down dispositivo} para desconectarme del dispositivo, y ya me quedo con la del móvil.
    Para configurarlo para que funcione con el móvil tendré que comprobar
    la ip de mi móvil buscando en la configuración.

    La cuenta de administrador de la instancia del móvil será la misma con la
    misma contraseña que la del ordenador.

    \section{Instalación actual}

    \subsection{Instalación en el móvil}
    Para instalarlo en el móvil es importante elegir en la configuración de la base de
    datos la base de datos SQLite, si no a mí actualmente no me funciona.

    \subsection{Instalación en el ordenador}
    Forgejo lo tengo instalado desde un contenedor en podman, que se inicia automáticamente como
    un servicio de systemd que se encuentra en \texttt{/etc/containers/systemd/forgejo.container}.
    Tal y como lo describe la documentación de Forgejo actualmente \href{https://forgejo.org/docs/latest/admin/installation/docker/}
    {aquí}, archivado \href{https://web.archive.org/web/20260307221713/https://forgejo.org/docs/latest/admin/installation/docker/}
    {aquí}. Y que actualmente tiene el siguiente contenido:

    \inputminted[breaklines, breakanywhere, breaksymbolleft="", breaksymbolright="", breakanywheresymbolpre="", breakanywheresymbolpost=""]
    {systemd}{inc/forgejo.container}

    Para que funcione son necesarios otros dos archivos en la misma carpeta, \texttt{/etc/containers/systemd/forgejo.network}
    y \texttt{/etc/containers/systemd/forgejo-data.volume}, que ambos tienen el siguiente contenido:

    \inputminted[breaklines, breakanywhere, breaksymbolleft="", breaksymbolright="", breakanywheresymbolpre="", breakanywheresymbolpost=""]
    {systemd}{inc/forgejo.network}

    \inputminted[breaklines, breakanywhere, breaksymbolleft="", breaksymbolright="", breakanywheresymbolpre="", breakanywheresymbolpost=""]
    {systemd}{inc/forgejo-data.volume}

    Ahora, Forgejo funciona con el usuario y la contraseña del administrados que tengo
    guardados tanto en papel como en mi gestor de contraseñas.

    El servidor se encuentra en \href{http://localhost:3000/}{http://localhost:3000/}.

    Actualmente el directorio en el que se encuentra todo lo relacionado con Forgejo es el siguiente:
    \texttt{/var/lib/containers/storage/volumes/forgejo-data/\_data/gitea/}, del que tendré que
    guardar copia de seguridad para no perder todos los repositorios.

    Para iniciar el servidor basta hacer \texttt{sudo systemctl start forgejo}.

    \section{Conectar desde las máquinas virtuales a Forgejo}
    Como se trata de una conexión desde la máquina virtual hacia el anfitrión, no
    es difícil de encontrar. Aunque no hay absolutamente nada de información sobre
    ello en el Internet. Un escaneo de las direcciones ip pasado a \texttt{nmap}, sobre
    el valor que devuelve \texttt{ip a}, que es de la forma 192.168.1.1/24, por ejemplo,
    devuelve en las máquinas virtuales que tengo que mi máquina anfitrión se encuentra en
    192.168.122.1. Así que allí puedo conectarme, a 192.168.122.1:3000 a través de las
    máquinas virtuales con http. El ssh no funciona porque parece que es importante
    en ese caso que el dominio de la web sea localhost.

    \section{Arreglar errores con ssh}
    Hay que tener cuidado, porque por defecto Forgejo tenía puesto como puerto de ssh el 222
    en el inicio de systemd, y aunque en los enlaces no especificara el
    puerto porque dentro del contenedor era el 22, fuera era el 222. Y por eso
    no funcionaba.

    \section{Configurar ssh}
    Es interesante la aproximación de Codeberg, que muestra en su guía sobre cómo crear
    las claves ssh, disponible \href{https://docs.codeberg.org/security/ssh-key/}{aquí},
    archivado \href{https://web.archive.org/web/20260205114835/https://docs.codeberg.org/security/ssh-key/}
    {aquí}, para que se puedan alternar distintas claves SSH para que se ponga
    automáticamente la de GitHub cuando me conecte allí, o a mi Instalación
    local cuando acceda a ella. El archivo de configuración donde se hace esto es
    en \texttt{~/.ssh/config}, que no expongo públicamente.

    \section{Copia de seguridad}
    \subsection{Forgejo del móvil}
    Todos los repositorios deben tener una instancia espejo establecida. Los pasos
    son los siguientes:

    \begin{enumerate}
        \item Crear un nuevo repositorio en el ordenador (siempre hay que seguir estos
        pasos al crearlo).
        \item Crear otro mismo repositorio con el mismo nombre que el de la instancia del móvil.
        \item Ir al apartado de configuración del repositorio en el ordenador.
        \item Ir a los «Mirror settings».
        \item Copiar el enlace que se me da en el repositorio del móvil http como si
        fuera a clonar el repositorio (del Forgejo del móvil).
        \item Pegar el enlace en el «Git remote repository URL».
        \item Activar la opción de sincronizar cuando se haga un «commit».
        \item Abrir el apartado de autorización, no usar la autorización por SSH, si no
        la de usuario y contraseña. Metiendo las mismas exactamente igual que en mi propia
        instancia.
        \item Una vez dado a agregar, darle al botón de sincronizar en el momento (si no
        se sincronizará más tarde).
    \end{enumerate}

    \subsection{Con Rclone}
    Se guarda automáticamente una copia del volumen exportado en un archivo \texttt{tar},
    actualmente se hace con el comando \texttt{sudo podman volume export forgejo-data -o forgejo.tar}.
    Esto lo hago con \texttt{crontab} con el superusuario, pues el contenedor necesita
    permisos elevados por el momento (lo podría cambiar). Para ello, hago \texttt{sudo crontab -e}
    y añado una línea como la que sigue:

    \texttt{@reboot podman volume export forgejo-data -o /home/usuario/Documentos/forgejo.tar}

    Como mi carpeta de documentos deberá estar sincronizada con Rclone en la nube, los
    datos estarán siempre en esa nube. Y en caso de pérdida solamente tengo que restaurar
    de la papelera de la cuenta y navegar con Rclone los archivos para buscar ese.
\end{document}
